\chapter{Matrice de couverture du cahier des charges initial}
\label{app:specification-coverage}

\newcommand{\covreal}{\textbf{\textcolor{ctgreen}{Réalisée}}}
\newcommand{\covadapt}{\textbf{\textcolor{ctblue}{Adaptée}}}
\newcommand{\covpartial}{\textbf{\textcolor{ctamber}{Partielle}}}
\newcommand{\covdeferred}{\textbf{\textcolor{ctgray}{Reportée}}}
\newcommand{\covvalidate}{\textbf{\textcolor{ctamber}{À valider}}}

\section{Objet et méthode}

Cette annexe confronte le cahier des charges fonctionnel version 1.0 transmis par
\CompanyName{} au produit disponible le 29 juillet 2026. Elle ne remplace pas ce
document initial et ne transforme pas une simulation en validation de
production. Les constats s'appuient sur le code source, le Product Backlog, la
spécification OpenAPI, le cahier de tests et les résultats automatisés.

\begin{longtable}{|C{2.4cm}|p{11.7cm}|}
  \caption{Légende de la matrice de couverture}
  \label{tab:coverage-legend}\\
  \hline
  \textbf{Statut} & \textbf{Interprétation}\\
  \hline
  \endfirsthead
  \hline
  \textbf{Statut} & \textbf{Interprétation}\\
  \hline
  \endhead
  \covreal & Le besoin principal est disponible de bout en bout dans le MVP et
  possède des preuves techniques adaptées.\\\hline
  \covadapt & L'objectif métier est couvert avec une solution ou une règle
  différente de l'option proposée initialement.\\\hline
  \covpartial & Une partie exploitable existe, mais au moins une capacité
  explicitement demandée n'est pas livrée.\\\hline
  \covdeferred & La fonctionnalité est volontairement sortie du MVP et conservée
  comme évolution.\\\hline
  \covvalidate & L'implémentation locale ne permet pas de prouver une cible de
  production ou un comportement matériel réel.\\
  \hline
\end{longtable}

\section{Modules fonctionnels}

\begingroup
\small
\begin{longtable}{|C{0.9cm}|p{2.7cm}|C{2.0cm}|p{8.5cm}|}
  \caption{Couverture des 18 modules fonctionnels initiaux}
  \label{tab:coverage-functional-modules}\\
  \hline
  \textbf{Réf.} & \textbf{Module} & \textbf{Statut} &
  \textbf{Couverture et écart constaté}\\
  \hline
  \endfirsthead
  \hline
  \textbf{Réf.} & \textbf{Module} & \textbf{Statut} &
  \textbf{Couverture et écart constaté}\\
  \hline
  \endhead
  M01 & Authentification & \covpartial &
  Connexion, déconnexion, inscription client, vérification email,
  réinitialisation, Google OAuth, invitations, RBAC et isolation sont réalisés.
  La double authentification n'est pas implémentée.\\\hline
  M02 & Gestion des bornes & \covpartial &
  Création, modification, suppression contrôlée, désactivation, géolocalisation,
  organisation, modèle, fabricant, puissance, identité OCPP, historique et états
  calculés sont disponibles. Le numéro de série et la version firmware ne sont
  pas encore conservés ; le pilotage du firmware relève de M13.\\\hline
  M03 & Connecteurs & \covreal &
  Une station possède plusieurs connecteurs numérotés avec prise, courant,
  puissance, état OCPP, erreur, source de disponibilité et QR d'entrée. Les
  changements de statut proviennent des événements de la borne.\\\hline
  M04 & Supervision temps réel & \covreal &
  Carte, listes, état calculé, occupation, connexion, dernier heartbeat, alertes,
  télémétrie disponible et actualisation Reverb sont intégrés. Une mesure absente
  n'est jamais remplacée par une valeur synthétique.\\\hline
  M05 & Sessions de recharge & \covadapt &
  Le parcours couvre préparation, préautorisation, démarrage OCPP, mesures,
  arrêt distant ou automatique, interruption, clôture et coût. La notion de
  pause n'est pas retenue pour une transaction OCPP 1.6J.\\\hline
  M06 & Utilisateurs & \covadapt &
  Super Admin, Administrateur, Opérateur, Technicien et Client remplacent les
  rôles ambigus Finance et Service Client. Les employés appartiennent à une
  organisation ; le client reste global et peut utiliser plusieurs réseaux.\\\hline
  M07 & RFID & \covpartial &
  Un identifiant OCPP virtuel condensé est associé au client et utilisé par
  \texttt{Authorize}. La gestion de badges physiques, leur lecture matérielle,
  rotation, blocage et expiration depuis l'interface restent à réaliser.\\\hline
  M08 & Véhicules & \covdeferred &
  Le profil véhicule a été retiré du MVP pour simplifier le parcours. La
  compatibilité est vérifiée par le type de connecteur choisi, sans conserver
  marque, modèle, immatriculation ou capacité de batterie.\\\hline
  M09 & Paiement & \covadapt &
  Le contrat \texttt{PaymentGateway} couvre autorisation, capture, libération,
  erreurs, webhooks, historique, reçus et factures. Carte, e-DINAR et D17 sont
  simulés avec WireMock ; paiement réel, remboursement et conformité bancaire
  sont reportés.\\\hline
  M10 & Tarification & \covpartial &
  Prix par kWh, frais de session, frais d'inactivité par minute, minimum,
  périodes de validité, affectations et remises d'abonnement sont disponibles.
  Les plages heures pleines/creuses et la tarification par puissance ou heure
  ne sont pas encore calculées.\\\hline
  M11 & Maintenance & \covreal &
  Plans préventifs ou correctifs, calendrier, récurrence, ordres de travail,
  affectation, SLA, états, pièces, commentaires, photos avant/après, historique
  et rapport final sont intégrés.\\\hline
  M12 & Alertes & \covpartial &
  Les défauts OCPP, pertes de communication, échéances et maintenances produisent
  alertes, notifications in-app et emails. Les signaux porte ouverte ou coupure
  dépendent d'une donnée de borne non encore testée ; SMS et push mobile ne sont
  pas livrés.\\\hline
  M13 & Firmware & \covdeferred &
  La plateforme conserve le modèle matériel et la version OCPP, mais pas de
  registre de versions firmware. Téléversement, signature, déploiement,
  progression, retour arrière et validation matérielle sont reportés
  conformément à US-38.\\\hline
  M14 & Rapports & \covpartial &
  Analyses par rôle, consommation, revenus, utilisation, disponibilité,
  maintenance, échanges internes et exports PDF, CSV et JSON sont réalisés.
  L'export Excel natif et les exports volumineux asynchrones restent à compléter.\\\hline
  M15 & Tableaux de bord & \covreal &
  Chaque rôle possède ses propres indicateurs, classements, tendances, alertes
  et cartes selon son périmètre. Les valeurs proviennent des services backend et
  exposent leur période de calcul.\\\hline
  M16 & Gestion documentaire & \covreal &
  Notices, certificats, contrats, garanties, photos et pièces de rapports peuvent
  être stockés de manière privée, prévisualisés, téléchargés et supprimés selon
  les permissions.\\\hline
  M17 & Paramétrage & \covpartial &
  Logo d'organisation, préférences email, fuseau personnel, rayon de recherche
  et contrôles globaux autorisés sont disponibles. TVA, SMS, internationalisation
  complète et paramètres de facturation de production restent incomplets.\\\hline
  M18 & Journal d'audit & \covreal &
  Les connexions et actions sensibles conservent acteur, événement, ressource,
  contexte et date. La consultation filtrée et les exports sont réservés au
  Super Admin.\\
  \hline
\end{longtable}
\endgroup

\section{Architecture et technologies}

\begin{longtable}{|C{1.0cm}|p{3.2cm}|C{2.0cm}|p{7.9cm}|}
  \caption{Couverture des orientations techniques}
  \label{tab:coverage-technical-choices}\\
  \hline
  \textbf{Réf.} & \textbf{Prévision} & \textbf{Statut} &
  \textbf{Décision réalisée}\\
  \hline
  \endfirsthead
  \hline
  \textbf{Réf.} & \textbf{Prévision} & \textbf{Statut} &
  \textbf{Décision réalisée}\\
  \hline
  \endhead
  T01 & Laravel 12 et PHP 8.3 & \covadapt &
  Laravel 13 est utilisé ; Composer accepte PHP à partir de 8.3 et
  l'environnement local fonctionne sous PHP 8.5.\\\hline
  T02 & API REST et JWT & \covadapt &
  L'API REST est réalisée et documentée en OpenAPI 3.1. La SPA utilise Sanctum,
  les cookies de session et CSRF plutôt qu'un JWT persistant.\\\hline
  T03 & Redis, queue et scheduler & \covadapt &
  Queue, cache et sessions utilisent PostgreSQL dans le MVP. Le scheduler et le
  worker sont configurés et exécutables ; Redis reste une option de déploiement,
  non une dépendance.\\\hline
  T04 & React, Ant Design, TypeScript, Axios et React Query & \covreal &
  Ces technologies structurent l'interface. React Hook Form, Zod et Zustand
  complètent les formulaires et états locaux.\\\hline
  T05 & ECharts & \covadapt &
  Recharts fournit les visualisations. Le changement de bibliothèque ne modifie
  pas les indicateurs attendus.\\\hline
  T06 & PostgreSQL ou MySQL & \covreal &
  PostgreSQL est le moteur retenu pour le développement et les requêtes métier.\\\hline
  T07 & WebSocket et Reverb ou Pusher & \covreal &
  Laravel Reverb et Echo diffusent les événements ; le protocole Pusher sert de
  transport compatible sans dépendre du service SaaS Pusher.\\\hline
  T08 & Google Maps ou OpenStreetMap & \covreal &
  Leaflet et OpenStreetMap affichent les cartes. Google Maps est uniquement
  utilisé comme destination externe optionnelle pour l'itinéraire.\\\hline
  T09 & OCPP 1.6 et 2.0.1 & \covpartial &
  OCPP 1.6 JSON est implémenté dans une passerelle Python séparée. OCPP 2.0.1
  reste une évolution.\\\hline
  T10 & Serveur OCPP dédié & \covreal &
  La passerelle WebSocket est indépendante de Laravel ; les événements sont
  transmis par API interne signée et les commandes sont réclamées par le gateway.\\
  \hline
\end{longtable}

\section{Messages OCPP}

\begin{longtable}{|p{4.2cm}|C{2.0cm}|p{8.0cm}|}
  \caption{Couverture du catalogue OCPP demandé}
  \label{tab:coverage-ocpp-messages}\\
  \hline
  \textbf{Message ou commande} & \textbf{Statut} & \textbf{Preuve fonctionnelle}\\
  \hline
  \endfirsthead
  \hline
  \textbf{Message ou commande} & \textbf{Statut} & \textbf{Preuve fonctionnelle}\\
  \hline
  \endhead
  \texttt{BootNotification} & \covreal &
  Enregistrement de la connexion et des informations de la borne.\\\hline
  \texttt{Heartbeat} & \covreal &
  Mise à jour de la dernière preuve de vie et détection planifiée des délais.\\\hline
  \texttt{StatusNotification} & \covreal &
  Projection des états et erreurs de chaque connecteur.\\\hline
  \texttt{Authorize} & \covreal &
  Validation d'un identifiant OCPP actif associé au client.\\\hline
  \texttt{StartTransaction} & \covreal &
  Création ou réconciliation de la session après confirmation de la borne.\\\hline
  \texttt{MeterValues} & \covreal &
  Énergie, puissance et coût provisoire issus des mesures reçues.\\\hline
  \texttt{StopTransaction} & \covreal &
  Clôture, compteur final et déclenchement de la capture du paiement.\\\hline
  \path{RemoteStartTransaction} & \covreal &
  Démarrage client guidé après détection du câble et préautorisation.\\\hline
  \path{RemoteStopTransaction} & \covreal &
  Arrêt manuel ou automatique demandé par la plateforme.\\\hline
  \texttt{Reset} & \covpartial &
  Le \emph{Soft Reset} est autorisé ; le \emph{Hard Reset} est refusé par mesure
  de sécurité dans le MVP.\\\hline
  \texttt{UnlockConnector} & \covreal &
  Commande distante tracée avec résultat retourné par la borne.\\\hline
  \texttt{ChangeAvailability} & \covreal &
  Extension utile ajoutée pour synchroniser le mode maintenance.\\\hline
  \texttt{ChangeConfiguration} & \covdeferred &
  Aucune modification distante arbitraire de la configuration n'est exposée.\\\hline
  \path{UpdateFirmware} et \path{FirmwareStatusNotification} & \covdeferred &
  Reportés avec le module firmware et la validation matérielle.\\\hline
  \path{GetDiagnostics} et \path{DiagnosticsStatusNotification} & \covdeferred &
  Les événements et journaux applicatifs existent, mais le transfert de
  diagnostics depuis une borne n'est pas implémenté.\\
  \hline
\end{longtable}

\section{Exigences non fonctionnelles}

\begin{longtable}{|C{1.0cm}|p{3.0cm}|C{2.0cm}|p{8.1cm}|}
  \caption{Couverture des exigences non fonctionnelles initiales}
  \label{tab:coverage-non-functional}\\
  \hline
  \textbf{Réf.} & \textbf{Exigence} & \textbf{Statut} &
  \textbf{Constat et limite}\\
  \hline
  \endfirsthead
  \hline
  \textbf{Réf.} & \textbf{Exigence} & \textbf{Statut} &
  \textbf{Constat et limite}\\
  \hline
  \endhead
  N01 & Page utilisable en moins de deux secondes & \covvalidate &
  Pagination, agrégats backend et tâches asynchrones sont en place, mais aucune
  campagne de performance représentative n'a encore validé ce seuil.\\\hline
  N02 & Disponibilité de 99,9\,\% & \covvalidate &
  La plateforme calcule la disponibilité des bornes ; la disponibilité du
  service lui-même devra être mesurée après déploiement et supervision.\\\hline
  N03 & HTTPS et chiffrement & \covvalidate &
  Les secrets sont exclus du dépôt, condensés ou chiffrés selon leur usage.
  TLS/WSS et la rotation de clés relèvent de l'environnement de production.\\\hline
  N04 & OAuth2, RBAC et audit & \covreal &
  Google OAuth, rôles, permissions, policies, isolation multi-tenant et journal
  d'audit sont appliqués côté serveur.\\\hline
  N05 & Interface responsive & \covpartial &
  Les parcours principaux sont adaptés aux ordinateurs et tablettes ; les
  parcours client essentiels fonctionnent sur mobile. Une recette responsive
  exhaustive reste à exécuter.\\\hline
  N06 & Chrome, Firefox et Edge & \covvalidate &
  Le socle utilise des API web standards, mais la matrice multi-navigateurs du
  cahier de tests doit être exécutée avant acceptation formelle.\\
  \hline
\end{longtable}

\section{Interfaces et livrables}

\begin{longtable}{|p{3.4cm}|C{2.0cm}|p{8.7cm}|}
  \caption{Couverture des interfaces attendues}
  \label{tab:coverage-interfaces}\\
  \hline
  \textbf{Interface} & \textbf{Statut} & \textbf{Implémentation}\\
  \hline
  \endfirsthead
  \hline
  \textbf{Interface} & \textbf{Statut} & \textbf{Implémentation}\\
  \hline
  \endhead
  Dashboard & \covreal & Tableaux de bord distincts pour les cinq rôles.\\\hline
  Carte GIS & \covreal & Carte Leaflet, filtres, géolocalisation et itinéraire.\\\hline
  Liste et détail des bornes & \covreal &
  Vues carte, cartes et tableau ; fiche avec connecteurs, supervision,
  télémétrie et sous-domaines.\\\hline
  Planning de maintenance & \covreal &
  Liste, calendrier, récurrence, affectation et cycle de vie.\\\hline
  Sessions et facturation & \covadapt &
  Sessions OCPP, paiements et documents fonctionnent avec un fournisseur simulé.\\\hline
  Rapports & \covreal &
  Analyses par rôle, rapports métier, échange interne et pièces jointes.\\\hline
  Paramètres & \covpartial &
  Préférences personnelles et contrôles globaux principaux disponibles ; certains
  paramètres de production restent à définir.\\
  \hline
\end{longtable}

\begin{longtable}{|p{4.2cm}|C{2.0cm}|p{7.9cm}|}
  \caption{État des livrables demandés}
  \label{tab:coverage-deliverables}\\
  \hline
  \textbf{Livrable} & \textbf{Statut} & \textbf{Emplacement ou observation}\\
  \hline
  \endfirsthead
  \hline
  \textbf{Livrable} & \textbf{Statut} & \textbf{Emplacement ou observation}\\
  \hline
  \endhead
  Cahier des charges & \covreal &
  Document fonctionnel version 1.0 transmis par l'entreprise et conservé comme
  référence initiale.\\\hline
  Maquettes UI & \covreal &
  Prototype navigable validé puis transposé dans l'application.\\\hline
  Architecture technique & \covreal &
  Chapitre 3, diagrammes et documentation du dépôt.\\\hline
  API REST & \covreal &
  \path{docs/api/openapi.json} et interface protégée \path{/docs/api}.\\\hline
  Code source & \covreal &
  Dépôt Git versionné avec frontend, backend, gateway et infrastructure locale.\\\hline
  Documentation & \covreal &
  Rapport, backlog, procédures OCPP, paiement et installation dans
  \path{docs/}.\\\hline
  Manuel utilisateur & \covreal &
  Source \LaTeX{} et PDF contrôlé dans \path{docs/manuals/user-manual}.\\\hline
  Manuel administrateur & \covreal &
  Source \LaTeX{} et PDF contrôlé dans \path{docs/manuals/admin-manual}.\\\hline
  Cahier de tests & \covreal &
  Source \LaTeX{}, résultats automatiques et fiches manuelles dans
  \path{docs/test-book}.\\\hline
  Procès-verbal de recette & \covdeferred &
  À préparer après exécution de la recette avec l'entreprise ; la décision et
  les réserves ne peuvent pas être anticipées par le stagiaire.\\
  \hline
\end{longtable}

\section{Lecture du résultat}

La couverture montre que le MVP respecte le coeur du besoin : administration
multi-organisation, supervision OCPP, disponibilité calculée, cartographie,
recharge, paiement simulé, exploitation et reporting. Les écarts majeurs ne sont
pas masqués : ils concernent les fonctions nécessitant du matériel, un contrat
externe, une infrastructure de production ou une extension fonctionnelle
distincte.

Cette matrice doit être mise à jour si une borne physique, un prestataire
financier, OCPP 2.0.1, les badges RFID ou le firmware entrent ultérieurement dans
le périmètre accepté.
